% !TeX program = lualatex
% Compile twice: lualatex open_problems_500.tex
% All references and prose are embedded; no BibTeX database or external inputs.
\documentclass[11pt,a4paper]{article}
\usepackage[margin=24mm,headheight=14pt]{geometry}
\usepackage{fontspec}
\setmainfont{Latin Modern Roman}
\setsansfont{Latin Modern Sans}
\setmonofont{Latin Modern Mono}
\usepackage{amsmath,amssymb,mathtools}
\usepackage{unicode-math}
\setmathfont{Latin Modern Math}
\usepackage{microtype}
\usepackage{enumitem,xurl,needspace}
\usepackage[dvipsnames]{xcolor}
\usepackage{hyperref}
\hypersetup{unicode=true,colorlinks=true,linkcolor=MidnightBlue,urlcolor=MidnightBlue,
 pdftitle={Research notebook: Problems 40--49},pdfauthor={Source-annotated compilation},bookmarksnumbered=true}
\usepackage{bookmark}
\usepackage{tocloft}
\setlength{\cftsecnumwidth}{3em}
\cftsetpnumwidth{2.5em}
\cftsetrmarg{4em}
\usepackage{titlesec}
\titleformat{\section}{\Large\bfseries\raggedright}{\thesection}{0.7em}{}
\usepackage{fancyhdr}
\pagestyle{fancy}\fancyhf{}
\fancyhead[L]{\small\sffamily Mathematical problem statements}
\fancyhead[R]{\small Source edition 22}
\fancyfoot[C]{\thepage}
\setlength{\parindent}{0pt}
\setlength{\parskip}{5pt plus 1pt}
\setlength{\emergencystretch}{3em}
\setcounter{tocdepth}{1}
\setcounter{secnumdepth}{1}
\setlist[itemize]{leftmargin=3em,itemsep=5pt,topsep=4pt,parsep=0pt}
\allowdisplaybreaks
\newcommand{\N}{\mathbb{N}}
\newcommand{\Z}{\mathbb{Z}}
\newcommand{\Q}{\mathbb{Q}}
\newcommand{\R}{\mathbb{R}}
\newcommand{\C}{\mathbb{C}}
\newcommand{\F}{\mathbb{F}}
\newcommand{\A}{\mathbb{A}}
\newcommand{\PP}{\mathbb P}
\newcommand{\E}{\mathbb{E}}
\newcommand{\eps}{\varepsilon}
\newcommand{\dd}{\,\mathrm d}
\newcommand{\sref}[2]{\hyperlink{source:#1:#2}{[S#2]}}
\newcommand{\eref}[2]{\hyperlink{extra:#1:#2}{[E#2]}}
% Turn the next switch off for a shorter reading copy. The quotations preserve
% every exactTarget field, including qualifications omitted from a short summary.
\newif\ifshowcatalogue
\showcataloguetrue
\newcommand{\cataloguescope}[1]{\ifshowcatalogue
 \paragraph{Exact scope in the supplied catalogue.}
 \begin{quote}\small #1\end{quote}\fi}

% Research additions dated 27 September 2026; compile twice with LuaLaTeX.
\numberwithin{equation}{section}
\begin{document}
\setcounter{section}{42}
% ================================================================
% problemId: problem.nonexistence-of-landau-siegel-zeros-for-quadratic-dirichlet-l-functions
% source releaseRank: 43; source releaseStatus: open
% exactTarget SHA-256: 766caffe5d517580dd01c8b519ab49351116fe349738b2bdf9f0319843418fbb
\clearpage
\section{\texorpdfstring{Nonexistence of Landau–Siegel zeros for quadratic Dirichlet L-functions}{Nonexistence of Landau–Siegel zeros for quadratic Dirichlet L-functions}}\label{43}
\subsection{Definitions and mathematical statement}
A primitive quadratic character $\chi$ has minimal modulus, its conductor $q$, and values in $\{0,\pm1\}$. Using its Dirichlet $L$-function, the uniform zero-free assertion is
\[
 \exists\delta>0\ \forall q\ge3\ \forall\chi\text{ primitive quadratic of conductor }q,
 \quad L(\beta,\chi)\ne0\quad\bigl(1-\delta/\log q<\beta<1\bigr).
\]
The same $\delta$ must work for all conductors, and only real $\beta$ are quantified here. A separate positive zero-free radius for each individual character would not establish the target.
\par\smallskip\textit{Formulation and concept references:} \sref{43}{1}.
\cataloguescope{Does there exist an absolute constant \ensuremath{\delta}>0 such that, for every primitive quadratic Dirichlet character \ensuremath{\chi} of conductor q\ensuremath{\ge}3, L(\ensuremath{\beta},\ensuremath{\chi})\ensuremath{\ne}0 for every real \ensuremath{\beta} with 1\ensuremath{-}\ensuremath{\delta}/log q<\ensuremath{\beta}<1? The logarithm is natural and the same \ensuremath{\delta} must work for all conductors.}
\subsection{Short English statement}
Can one uniformly exclude exceptionally close real zeros of quadratic Dirichlet $L$-functions near one?
% BEGIN RESEARCH ADDITION 43
\subsection{Research attempt}
\paragraph{Current frontier.}
Lu--Zaman--Zhao report a rigorous computational exclusion of real zeros in
$\sigma\ge1-1/(5\log q)$ for quadratic characters with $q\le10^{10}$
\sref{43}{1}, Theorem 1.1. Their verified finite range is not an
all-conductor theorem. Their algorithm's asymptotic termination analysis
uses GRH, separately from the finite verifications. The established
zero-free-region theorem allows at most one exceptional real zero rather
than excluding that zero. The distinction between lower bounds for
$L(1,\chi)$ and a zero-free interval of width $1/\log q$ is analyzed by
Friedlander--Iwaniec \eref{43}{1}, Sections 1--2. These are the relevant
inputs; a general proof claim retained in [S2] is not used as a theorem.

\paragraph{Attack route.}
Three possible routes are: control the variation of $L$ relative to its
value at one; exploit positivity of the coefficients of $\zeta L$; or
obtain a uniform lower bound for the positive prime sums used by the
numerical exclusion method. The first requires an estimate proportional
to $L(1,\chi)$, not just an absolute derivative estimate. The second must
continue through a pole without illegally continuing a positive series.
The third requires uniform information on splitting primes in a range
where exceptional zeros are precisely an obstruction. We carry out the
first route explicitly and test the positivity shortcut against a model.

\paragraph{Attempt: an unconditional derivative calculation.}
Let $q\ge3$, $\ell=\log q>1$, and let $\chi$ be primitive quadratic.
It is nonprincipal. Set $A(t)=\sum_{n\le t}\chi(n)$. Periodicity and the
zero sum of a nonprincipal character over one full period give
$A(q)=0$ and $|A(t)|\le q$. Partial summation yields, for real $s>0$,
\[
 L(s,\chi)=\sum_{n\le q}\frac{\chi(n)}{n^s}
               +s\int_q^\infty A(t)t^{-s-1}\,dt. \tag{43.1}
\]
The integral converges absolutely and agrees with the continuation from
$s>1$. On any compact subset of $s>0$, its derivative is dominated by an
integrable multiple of $(1+\log t)t^{-1-\sigma_0}$, where
$\sigma_0>0$. Differentiation is therefore justified, giving
\begin{align*}
 L'(s,\chi)={}&-\sum_{n\le q}\frac{\chi(n)\log n}{n^s}
       +\int_q^\infty A(t)t^{-s-1}\,dt\\
       &-s\int_q^\infty A(t)(\log t)t^{-s-1}\,dt. \tag{43.2}
\end{align*}
Take $1-1/(4\ell)\le s\le1$. Then $s\ge3/4$ and
$q^{1-s}\le e^{1/4}<2$. The finite sum has absolute value at most
\[
 e^{1/4}\sum_{n\le q}\frac{\log n}{n}
 \le e^{1/4}\ell(1+\ell)<4\ell^2.
\]
The two integral terms in (43.2) together have absolute value at most
\[
 q^{1-s}\left(\frac1s+\log q+\frac1s\right)
 \le2\left(\ell+\frac83\right)<8\ell^2.
\]
In particular the deliberately nonoptimal bound
\[
 \sup_{1-1/(4\ell)\le s\le1}|L'(s,\chi)|\le16\ell^2 \tag{43.3}
\]
is uniform in $q$ and $\chi$. The order of this estimate is familiar from
\eref{43}{1}, equation (1.4); the calculation above supplies an elementary
proof, not an improvement on that literature.

If a real zero $\beta$ lies in this interval, the fundamental theorem of
calculus and the classical fact $L(1,\chi)>0$ imply
\[
 0<L(1,\chi)=\int_\beta^1 L'(s,\chi)\,ds
       \le16(1-\beta)\ell^2. \tag{43.4}
\]
Here positivity at one is the usual nonvanishing and positivity theorem
for real nonprincipal Dirichlet characters; it is also implicit in the
special-value bounds discussed in \sref{43}{1} and \eref{43}{1}.
Thus even a hypothetical uniform estimate
$L(1,\chi)\ge c/\ell$ would yield from this route only
$1-\beta\ge c/(16\ell^3)$ inside the working interval. That is two powers
of $\log q$ short of the required scale. Increasing the finite verification
range does not repair this asymptotic loss.

\paragraph{Attempt: a relative-variation criterion.}
Here is a sufficient hypothesis that identifies the missing cancellation
without assuming a value for $L(1,\chi)$:
\[
 \exists C>0\ \ \forall\chi,q:\quad
 \sup_{1-1/(4\log q)\le s\le1}|L'(s,\chi)|
          \le C(\log q)L(1,\chi). \tag{43.5}
\]
If (43.5) holds, the target follows with
$\delta=\min\{1/4,1/(2C)\}$. Indeed, a zero with
$0<1-\beta<\delta/\log q$ would give
\[
 L(1,\chi)\le C(1-\beta)(\log q)L(1,\chi)
                  <\tfrac12 L(1,\chi),
\]
a contradiction. This is a complete proof of the conditional implication,
not a proof of (43.5). A one-sided integral estimate could be enough;
absolute control over the whole interval is stronger than what the
contradiction uses. In particular (43.5) should not be described as an
established equivalence or as an independently easier conjecture.

To seek (43.5), separate the principal part of (43.2) into sign classes.
Its finite sum is
\[
 -\sum_{\substack{n\le q\\\chi(n)=1}}\frac{\log n}{n^s}
 +\sum_{\substack{n\le q\\\chi(n)=-1}}\frac{\log n}{n^s}.
\]
Bounding both sums separately destroys the cancellation that (43.5)
requires when $L(1,\chi)$ is small. Character periodicity supplies (43.3)
but does not relate this cancellation to the small special value.
Friedlander--Iwaniec's more refined relation to a smoothed positive
coefficient sum explains why that link is substantive
\eref{43}{1}, Proposition 2.1. No such uniform relative estimate has been
proved in this attempt.

\paragraph{Attempt: testing positivity across the pole.}
For $s>1$ the product has the absolutely convergent expansion
\[
 Z(s)=\zeta(s)L(s,\chi)=\sum_{n\ge1}\frac{a(n)}{n^s},
 \qquad a(n)=\sum_{d\mid n}\chi(d)\ge0. \tag{43.6}
\]
The last assertion follows prime by prime. If $\chi(p)=1$, then
$a(p^j)=j+1$; if $\chi(p)=-1$, it equals one for even $j$ and zero for
odd $j$; if $\chi(p)=0$, it equals one. Multiplicativity proves the
claim, including $a(m^2)\ge1$. Hence $Z(s)>0$ for $s>1$ and its
alternating real derivatives have the signs supplied by a positive
Laplace transform. This calculation is checked on small conductors in
the accompanying exact script.

An attempted continuation of these signs to $s<1$ is invalid. To expose
the exact logical failure, put $t=s-1$ and, for any $\epsilon,C>0$, define
\[
 Z_\epsilon(t)=C\left(1+\frac{\epsilon}{t}\right)
      =C\int_{[0,\infty)}e^{-tu}
                          (\delta_0+\epsilon\,du),\qquad t>0.
 \tag{43.7}
\]
It is positive and completely monotone for $t>0$, has a positive-residue
simple pole at zero, and nevertheless its meromorphic continuation has a
zero at $t=-\epsilon$. After removing the pole,
$L_\epsilon(t)=tZ_\epsilon(t)=C(t+\epsilon)$ has
$L_\epsilon'(0)/L_\epsilon(0)=1/\epsilon$. Thus positivity of the transform
alone cannot control the relative variation. This model is neither a
Dirichlet $L$-function nor a counterexample to the problem: it lacks the
arithmetic Euler factors and functional equation. It falsifies only the
proposed inference from positivity and a simple pole.

\paragraph{Stress test.}
The finite computations in the cited paper prove a finite statement, not
termination of a conductor-by-conductor argument for all $q$. Applying its
GRH-conditional termination estimate here would assume a substantially
stronger unresolved zero-location statement. Likewise, excluding all but
one exceptional character in a range does not exclude that last character.
The derivative calculation covers the stipulated real interval and uses
uniform constants; it does not claim complex zero-freeness. Every exchange
of derivative and integral in (43.2) was made in $s>0$ with an explicit
integrable majorant, whereas no use of the series (43.6) was made left of
its convergence domain.

For any fixed finite collection of characters, positivity and continuity
at one give a positive minimum zero-free radius after rescaling by
$\log q$. That elementary compactness argument applies to a finite set,
not to an unbounded sequence of conductors. The potential counterexample
regime is precisely a sequence for which the rescaled radii tend to zero.
None of the inequalities derived above excludes that regime.

\paragraph{Outcome.}
\textbf{Outcome: Conditional reduction.}
The target follows from the explicit relative-variation estimate (43.5).
The supplied proof of the implication handles all conductors uniformly.
What was proved unconditionally is a coarse derivative bound and the
resulting weaker special-value-to-zero-distance estimate, together with a
countermodel to a positivity-only shortcut. These auxiliary calculations
are not claimed to be historically new. The missing cancellation in
(43.5), or an adequate one-sided substitute, has not been established;
there is no unconditional exclusion of all Landau--Siegel zeros here.
% END RESEARCH ADDITION 43
\subsection{Sources}
\begin{itemize}\raggedright
\item[\textnormal{[S1]}]\hypertarget{source:43:1}{} Rick F. Lu, Asif Zaman and Haonan Zhao, Numerical Computations Concerning Landau–Siegel Zeros; arXiv2602.03626v1 February3,2026.
\url{https://arxiv.org/pdf/2602.03626v1}.
{\small\textit{Catalogue formulation source.}}
\item[\textnormal{[S2]}]\hypertarget{source:43:2}{} On the Extended, Generalized, and Grand Riemann Hypotheses: A Unified Approach via General Properties of L-Functions, v18.
\url{https://www.preprints.org/manuscript/202506.0481/v18}.
{\small\textit{Catalogue research/status reference; not a proof endorsement.}}
\item[\textnormal{[S3]}]\hypertarget{source:43:3}{} Numerical Computations concerning Landau–Siegel Zeros.
\url{https://arxiv.org/html/2602.03626v1}.
{\small\textit{Catalogue research/status reference; not a proof endorsement.}}
\item[\textnormal{[S4]}]\hypertarget{source:43:4}{} Landau-Siegel zeros of Rankin-Selberg L-functions.
\url{https://arxiv.org/abs/2601.04189}.
{\small\textit{Catalogue research/status reference; not a proof endorsement.}}
% BEGIN RESEARCH SOURCES 43
\item[\textnormal{[E1]}]\hypertarget{extra:43:1}{} John B. Friedlander and Henryk Iwaniec,
\emph{A note on Dirichlet $L$-functions}, Expositiones Mathematicae 36 (2018),
343--350; arXiv:1701.03771v1 (2017), Section 1 and Proposition 2.1.
\url{https://arxiv.org/html/1701.03771v1}.
{\small\textit{The distinction between special-value estimates and real zero-free intervals. Consulted 27 September 2026.}}
% END RESEARCH SOURCES 43
\end{itemize}

\end{document}
